\chapter{配合物的反应机理}\label{ch:b3:complex-mechanisms}

将铬(III)盐溶于用氧-18 标记的水中，
需要数天，标记的水才能取代与金属结合的水分子；
而对于铜(II)盐，交换在混合所需的时间内即告完成。
两者的反应本质相同：一个水分子离开金属离子，
另一个水分子随即占据其位置，但速率因离子而差异极大，
其中大部分差异可由 $d$ 电子数预测。
配体如何被取代，决定了铂类抗癌药物如何制备、如何发挥作用；
电子如何在金属离子间移动，则决定了呼吸作用和每一种电池的速率。
本章讨论配合物的两类反应：取代与电子转移，
其中电子转移借助
鲁道夫·马库斯的理论进行讨论。

\begin{recall}
本科二年级分册介绍了配合物、配体的齿数、桥联配体、
$面式$ 与 $经式$ 异构体、作为基元反应的配体交换，
以及高自旋和低自旋构型下的晶体场稳定化能；
本科一年级分册讲述了反应机理，以及稳态近似与决速步
近似。\Cref{ch:b3:rate-theories}给出了艾林方程，
并给出 $\Delta^{\ddagger}S^\circ$ 的含义，\cref{ch:b3:electrode-kinetics} 讨论
电极上的电子转移，以及\cref{ch:b3:complex-spectra-magnetism}
讨论配体场谱项和扬–泰勒效应。
\end{recall}

\section{活性与惰性}

\begin{definition}[活性配合物与惰性配合物]\label{def:b3:complex-mechanisms:labile}
若配合物的配体被迅速取代，则称其为\emph{活性配合物}\index{活性配合物}，
配体被缓慢取代的，称为\emph{惰性配合物}\index{惰性配合物}；
按陶贝的定义，配合物称为活性的，是指它
与配体在 \qty{0.1}{mol/L} 和 \qty{25}{\celsius} 下的反应在约一
分钟内完成。活性是一种动力学性质，与配合物的热力学稳定性
无关。
\end{definition}

水交换 $\mathrm{[M(H_2O)_6]^{n+} + H_2O^{\ast} \to [M(H_2O)_5(H_2O^{\ast})]^{n+} + H_2O}$ 是最简单的取代反应，
并确立了一个标度。主族离子在
半径较大、电荷较少时交换较快。在过渡金属离子中，具有 $d^3$
（\ce{Cr^{3+}}）或低自旋 $d^6$（\ce{Co^{3+}}、\ce{Rh^{3+}}、\ce{Ir^{3+}}）构型的离子呈惰性，大多数
第二、三周期金属也是如此，而 $d^9$（\ce{Cu^{2+}}）和高自旋 $d^4$（\ce{Cr^{2+}}）的离子
因扬–泰勒效应而发生畸变，活性极高。

\begin{proposition}[配体场活化能]\label{prop:b3:complex-mechanisms:lfae}
若取代反应经过五配位的四方锥过渡态，
则进入该过渡态时配体场稳定化能的损失（即配体场活化能）
在最简单的仅含 $\sigma$ 作用的模型中，对低自旋 $d^6$ 最大（$4\,Dq$），
其次为 $d^3$ 和 $d^8$（$2\,Dq$）；对 $d^0$、高自旋 $d^5$ 和 $d^{10}$ 为零；
对高自旋 $d^4$ 和 $d^9$ 则为负。
\end{proposition}

\begin{proof}
将每个配体看作使 $d$ 轨道升高 $e_\sigma$ 乘以该轨道在配体方向上的
角向振幅的平方，并在瓣轴上归一化为 1：轴向配体
使 $d_{z^2}$ 得 1，赤道向配体使 $d_{z^2}$ 得 $\frac14$、$d_{x^2-y^2}$ 得 $\frac34$；$t_{2g}$ 轨道指向配体之间，
不受影响。八面体中，$d_{z^2}$ 与 $d_{x^2-y^2}$ 各得 $3e_\sigma$
（$\Delta_{\mathrm o} = 3e_\sigma = 10Dq$）；去掉一个轴向配体的四方锥中，$d_{z^2}$ 得 $2e_\sigma$，
$d_{x^2-y^2}$ 仍为 $3e_\sigma$。$n$ 个电子的稳定化能等于其能量减去 $n/5$ 乘以
五个能级之和（即重心）。低自旋 $d^6$（$t_{2g}^6$）的
八面体中为 $0 - \frac65 \times 6e_\sigma$，四方锥中为 $0 - \frac65 \times 5e_\sigma$，损失为 $1.2e_\sigma = 4Dq$。对 $d^3$：$\frac35e_\sigma =
2Dq$；对 $d^8$（$t_{2g}^6e_g^2$）：$(5 - \frac85 \times 5) - (6 - \frac85 \times 6) = 0.6e_\sigma$。对 $d^9$：$(7 - 9) - (9 -
10.8) = -0.2e_\sigma$.
\end{proof}

\begin{omfigure}
\begin{tikzpicture}
\begin{axis}[width=0.72\linewidth, height=5.4cm, xmin=-0.6, xmax=10.6, ymin=-1.2, ymax=4.6,
  axis x line=bottom, axis y line=left, ycomb, xtick={0,1,2,3,4,5,6,7,8,9,10},
  xlabel={$n$ ($d^n$)}, ylabel={活化能 / $Dq$}, tick label style={font=\small},
  label style={font=\small}]
\draw[black!40] (axis cs:-0.6,0) -- (axis cs:10.6,0);
\addplot[omDef, line width=4pt, xshift=-2.5pt] table[x=n, y=hs] {figdata/out/bachelor-3-ligand-field-activation.dat};
\addplot[omThm, line width=4pt, xshift=2.5pt, unbounded coords=discard] table[x=n, y=ls] {figdata/out/bachelor-3-ligand-field-activation.dat};
\end{axis}
\end{tikzpicture}

{\small 解离途径经由四方锥时的配体场活化能，取自命题中的
仅含 $\sigma$ 的模型（蓝色：高自旋；
红色：低自旋）。最大值
出现在低自旋 $d^6$、$d^3$ 和 $d^8$ 处，与惰性或反应缓慢的离子（\ce{Co^{3+}}、\ce{Cr^{3+}}、\ce{Ni^{2+}}）相对应；
负值出现在 $d^4$ 和 $d^9$ 处，对应活性极高的扬–泰勒离子。}
\end{omfigure}

\begin{definition}[活化体积]\label{def:b3:complex-mechanisms:activation-volume}
称基元反应的\emph{活化体积}\index{活化体积} $\Delta V^{\ddagger}$ 为
活化络合物的偏摩尔体积
与反应物的偏摩尔体积之差。
\end{definition}

\begin{proposition}[速率常数的压强依赖性]\label{prop:b3:complex-mechanisms:pressure}
恒温下，$\displaystyle\Bigl(\frac{\partial\ln k}{\partial p}\Bigr)_T = -\frac{\Delta V^{\ddagger}}{RT}$。
\end{proposition}

\begin{proof}
由艾林方程 $\ln k = \text{常数} - \Delta^{\ddagger}G/RT$，并由 $(\partial\Delta^{\ddagger}G/\partial p)_T = \Delta V^{\ddagger}$，
与任一吉布斯自由能的情形相同。
\end{proof}

配体离去的步骤有 $\Delta V^{\ddagger} > 0$；配体进入的步骤有 $\Delta V^{\ddagger}
< 0$。在数百兆帕下测得的活化体积，是判断取代机理
最直接的依据，它较少受溶剂化的干扰，
较之活化熵更为可靠。

\section{取代机理}

\begin{definition}[解离机理、缔合机理与交换机理]\label{def:b3:complex-mechanisms:mechanisms}
在\emph{解离机理}\index{解离机理}（D）中，离去的
配体先离开，生成配位数较低的中间体；在
\emph{缔合机理}\index{缔合机理}（A）中，进入的配体
先结合，生成配位数较高的中间体；在
\emph{交换机理}\index{交换机理}（I）中，两个配体
在一步内完成交换，具有键断裂（$I_d$）或键生成（$I_a$）的
特征。
\end{definition}

\begin{proposition}[D 机理与 A 机理的速率方程]\label{prop:b3:complex-mechanisms:d-a-laws}
对于 $\mathrm{ML_5X + Y \to ML_5Y + X}$：D 机理（$\mathrm{ML_5X} \rightleftharpoons \mathrm{ML_5} + \mathrm X$，$k_1$，$k_{-1}$；$\mathrm{ML_5} + \mathrm Y
\to \mathrm{ML_5Y}$，$k_2$）给出
\[
  v = \frac{k_1k_2[\mathrm{ML_5X}][\mathrm Y]}{k_{-1}[\mathrm X] + k_2[\mathrm Y]},
\]
对配合物为一级，且在 $[\mathrm Y]$ 很大时与 $[\mathrm Y]$ 无关；A 机理
经一个处于稳态的七配位中间体，给出 $v =
k[\mathrm{ML_5X}][\mathrm Y]$，对每一反应物各为一级。
\end{proposition}

\begin{proof}
D：对 $\mathrm{ML_5}$ 取稳态，$k_1[\mathrm{ML_5X}] = (k_{-1}[\mathrm X] + k_2[\mathrm Y])[\mathrm{ML_5}]$，且 $v = k_2[\mathrm{ML_5}][\mathrm Y]$。
当 $[\mathrm Y]$ 很大时，$v \to k_1[\mathrm{ML_5X}]$。A：$\mathrm{ML_5X} + \mathrm Y \rightleftharpoons \mathrm{ML_5XY}$（$k_a$，$k_{-a}$），$\mathrm{ML_5XY} \to
\mathrm{ML_5Y} + \mathrm X$（$k_b$）；由稳态可得 $v = \frac{k_ak_b}{k_{-a} + k_b}[\mathrm{ML_5X}][\mathrm Y]$。
\end{proof}

在水中，进入的配体通常远低于
溶剂的浓度，真正的第一步是配合物与
配体相遇。

\begin{definition}[艾根–威尔金斯机理]\label{def:b3:complex-mechanisms:eigen-wilkins}
称水合配合物取代的\emph{艾根–威尔金斯机理}\index{艾根–威尔金斯机理}为
一个快速预平衡，生成
\emph{外层配合物}\index{外层配合物}，其中进入的配体紧挨配合物、
位于其配位层之外，
随后的决速步是该配体与一个水分子的交换。
\end{definition}

\begin{theorem}[艾根–威尔金斯速率方程]\label{thm:b3:complex-mechanisms:eigen-wilkins}
设外层平衡常数为 $K_{\mathrm{os}}$，交换速率常数为
$k_i$，则在 $[\mathrm Y]$ 过量时，观测到的准一级速率常数为
\[
  k_{\mathrm{obs}} = \frac{K_{\mathrm{os}}k_i[\mathrm Y]}{1 + K_{\mathrm{os}}[\mathrm Y]} .
\]
\end{theorem}

\begin{proof}
配合物在游离形式与外层形式之间分配，$[\mathrm{M\cdot Y}] = K_{\mathrm{os}}[\mathrm M][\mathrm Y]$，
总浓度为 $[\mathrm M]_{\mathrm{tot}} = [\mathrm M](1 + K_{\mathrm{os}}[\mathrm Y])$；速率为 $k_i[\mathrm{M\cdot Y}] = k_iK_{\mathrm{os}}[\mathrm Y][\mathrm M]_{\mathrm{tot}}/(1 + K_{\mathrm{os}}[\mathrm Y])$。
\end{proof}

当 $[\mathrm Y]$ 很低时，反应为二级，$k = K_{\mathrm{os}}k_i$；$K_{\mathrm{os}}$ 可由电荷与最近接近距离估算，
对于同电荷的所有配体几乎相同，
因此 $k_i$（即水交换步骤）控制
速率：水合离子上的取代几乎与进入的配体无关，
这正是解离型交换机理的特征。

\begin{definition}[共轭碱机理]\label{def:b3:complex-mechanisms:conjugate-base}
称氨配合物碱水解的\emph{共轭碱机理}\index{共轭碱机理}为
（如 $\mathrm{[Co(NH_3)_5X]^{2+} + OH^- \to [Co(NH_3)_5(OH)]^{2+} + X^-}$ 所示）：在快速预平衡中，氨配体先脱质子，
随后，氨基共轭碱发生解离
失去 $\mathrm X^-$，其 $\mathrm{NH_2^-}$ 配体
使配合物成为活性配合物。
\end{definition}

\begin{proposition}[碱水解的速率方程]\label{prop:b3:complex-mechanisms:conjugate-base}
设脱质子平衡常数为 $K$，
共轭碱的解离速率常数为 $k$，则 $v = \dfrac{kK[\text{配合物}]_{\mathrm{tot}}[\ce{OH-}]}{1 + K[\ce{OH-}]}$，在 $[\ce{OH-}]$ 很低时为二级。
\end{proposition}

\begin{proof}
与艾根–威尔金斯速率方程的证明相同：配合物在其酸式与
共轭碱式之间按 $1 : K[\ce{OH-}]$ 的比例分配，且只有后者反应。
\end{proof}

仅凭速率方程，无法将该机理与
\ce{OH-}的直接进攻区分开来；证据是碱水解需要一个 N--H 质子（不含此质子的配合物
不表现此现象），且氨配体上的质子与溶剂
交换的速率远快于水解。

\begin{method}[取代机理的判别]\label{met:b3:complex-mechanisms:diagnose}
\begin{enumerate}
  \item 对进入基团的依赖：速率与 Y 无关（或趋于饱和），
        为解离型；速率与 $[\mathrm Y]$ 成正比且对 Y 的性质敏感，
        为缔合型。
  \item 活化体积：为正，为解离型；为负，为缔合型。
  \item 活化熵：对 D 为正，对 A 为负（需谨慎，因为
        溶剂化也有贡献）。
  \item 对离去基团的依赖：速率随 M--X 键的强度变化，
        表明过渡态中存在键断裂。
\end{enumerate}
\end{method}

\begin{omfigure}
\begin{tikzpicture}
\begin{axis}[name=pd, omprofile, width=0.46\linewidth, height=4.6cm, xmin=0, xmax=10, ymin=0, ymax=10,
  xlabel={反应坐标}, ylabel={能量}, title={\small D：$\mathrm{ML_5X} \to \mathrm{ML_5} + \mathrm X \to \mathrm{ML_5Y}$}]
\addplot[omDef, very thick, smooth] coordinates {(0,2) (1.5,2.1) (3,7.5) (4.2,5.6) (5.0,5.4) (5.8,5.6) (7,7.0) (8.5,1.6) (10,1.5)};
\node[font=\scriptsize, anchor=north, align=center] at (axis cs:5,5.2) {ML$_5$（配位数 5）};
\end{axis}
\begin{axis}[at={(pd.east)}, anchor=west, xshift=1.0cm, omprofile, width=0.46\linewidth, height=4.6cm,
  xmin=0, xmax=10, ymin=0, ymax=10, xlabel={反应坐标}, ylabel={能量},
  title={\small A：$\mathrm{ML_5X} + \mathrm Y \to \mathrm{ML_5XY} \to \mathrm{ML_5Y}$}]
\addplot[omThm, very thick, smooth] coordinates {(0,2) (1.5,2.1) (3,6.5) (4.2,4.4) (5.0,4.2) (5.8,4.4) (7,7.2) (8.5,1.6) (10,1.5)};
\node[font=\scriptsize, anchor=north, align=center] at (axis cs:5,4.0) {ML$_5$XY（配位数 7）};
\end{axis}
\end{tikzpicture}

{\small 解离型与缔合型取代的能量曲线
（示意）：二者都经过一个中间体，D 中中间体的配位数（CN）较低，
A 中则较高。在交换机理中
势阱消失，只剩下一个过渡态。}
\end{omfigure}

经四方锥中间体发生取代的八面体配合物，通常
保持其构型（$cis$ 仍为 $cis$）；而经三角双锥中间体
取代的配合物（许多钴(III)配合物在碱水解中即如此）
可得 $cis$ 与 $trans$ 产物的混合物。

\section{平面四方形取代与反位效应}

$d^8$ 离子的平面四方形配合物（\ce{Pt^{2+}}、\ce{Pd^{2+}}、\ce{Au^{3+}}、与强配体结合的 \ce{Ni^{2+}}）
配位不饱和：进入的配体可沿
空的轴向接近。其取代属缔合型，
经五配位三角双锥过渡态进行。

\begin{proposition}[双项速率方程]\label{prop:b3:complex-mechanisms:square-planar}
在配位溶剂 S 中，取代反应 $\mathrm{[PtL_3X] + Y \to [PtL_3Y] + X}$ 遵循
$k_{\mathrm{obs}} = k_1 + k_2[\mathrm Y]$（Y 过量时）。
\end{proposition}

\begin{proof}
两条平行的缔合路径：Y 直接进攻（$k_2[\mathrm Y]$）以及
溶剂进攻（溶剂大量过量，为准一级 $k_1$），生成 $[\mathrm{PtL_3S}]$，
随即被 Y 迅速转化。平行的一级路径速率相加。
\end{proof}

\begin{definition}[反位效应，反位影响]\label{def:b3:complex-mechanisms:trans-effect}
\emph{反位效应}\index{反位效应}是指一个配体被其反位配体取代时的加速作用，
是作用于过渡态的动力学效应。
\emph{反位影响}\index{反位影响}是指基态中与某配体处于反位的键
被削弱，可由键长和振动频率看出，
是一种热力学效应。
\end{definition}

反位效应的顺序为 $\ce{CO}$、$\ce{CN-}$、$\ce{C2H4} > \ce{PR3}$、$\ce{H-} > \ce{CH3-} > \ce{I-}$、$\ce{SCN-} >
\ce{Br-} > \ce{Cl-} > \ce{NH3}$、$\ce{H2O}$。强 $\sigma$ 给体削弱反位键（反位影响）；
强 $\pi$ 受体稳定五配位过渡态，其方式是
从金属夺取电子密度。二者均使反应速率增大。

\begin{method}[规划铂(II)异构体的合成]\label{met:b3:complex-mechanisms:pt-synthesis}
\begin{enumerate}
  \item 每一步中，被替换的配体是与现存的反位效应最强的
        配体处于反位的那一个。
  \item 安排各步加入配体的顺序，使该规则导出所需的异构体。
\end{enumerate}
\end{method}

\begin{omfigure}
\setchemfig{atom sep=2.2em}
\schemestart
\chemfig{Cl-Pt(-[2]Cl)(-[6]Cl)-Cl}
\arrow{->[\ce{NH3}]}
\chemfig{Cl-Pt(-[2]Cl)(-[6]Cl)-NH_3}
\arrow{->[\ce{NH3}]}
\chemfig{Cl-Pt(-[2]NH_3)(-[6]Cl)-NH_3}
\schemestop

\bigskip
\schemestart
\chemfig{H_3N-Pt(-[2]NH_3)(-[6]NH_3)-NH_3}
\arrow{->[\ce{Cl-}]}
\chemfig{H_3N-Pt(-[2]NH_3)(-[6]NH_3)-Cl}
\arrow{->[\ce{Cl-}]}
\chemfig{Cl-Pt(-[2]NH_3)(-[6]NH_3)-Cl}
\schemestop

{\small 反位效应实例（略去电荷）。上：自 $\ce{[PtCl4]^{2-}}$ 出发，第二个氨分子
取代与另一氯离子呈反位的氯离子（Cl$^-$ 的反位效应更大），
得 $cis$ 异构体，即顺铂。下：自 $\ce{[Pt(NH3)4]^{2+}}$ 出发，第二个氯离子
取代与第一个氯离子呈反位的氨分子，得 $trans$ 异构体。}
\end{omfigure}

\section{电子转移}

\begin{definition}[外层与内层电子转移]\label{def:b3:complex-mechanisms:electron-transfer}
在\emph{外层电子转移}\index{外层电子转移}中，
电子在两个配位层保持完整的配合物之间传递；
在\emph{内层电子转移}\index{内层电子转移}中，电子
经由一个桥联两个金属的配体传递，该配体处于
瞬时双核配合物中。\emph{自交换反应}\index{自交换反应}是指
同一电对两种氧化态之间的电子转移，且无
净化学变化，例如 $\ce{[Fe(H2O)6]^{3+} + [Fe(H2O)6]^{2+}}$。
\end{definition}

亨利·陶布于 1953 年以一个经精心选择的反应证明了内层途径，所选反应
使产物本身即能说明问题。钴(III) 为惰性，钴(II) 为活性；铬(II)
为活性，铬(III) 为惰性。当 $\ce{[Co(NH3)5Cl]^{2+}}$ 被 $\ce{[Cr(H2O)6]^{2+}}$ 还原时，
在酸中，所生成的铬(III) 均携带氯离子，以 $\ce{[Cr(H2O)5Cl]^{2+}}$ 的形式存在，即使
游离氯离子存在时亦然：该氯离子必定桥联了两个金属，
在电子转移的瞬间即与惰性的铬(III) 成键，并
随后被活性的钴(II) 释放。

\begin{omfigure}
\begin{tikzpicture}[font=\small]
  \begin{scope}[scale=0.8]
    \pic[transform shape] (a) at (0,0) {omoctahedron};
    \pic[transform shape] (b) at (2.0,0) {omoctahedron};
  \end{scope}
  \node[anchor=north east] at (0.3,-1.0) {$\mathrm{Co^{III}(NH_3)_5}$};
  \node[anchor=north west] at (1.3,-1.0) {$\mathrm{Cr^{II}(H_2O)_5}$};
  \node[omThm, anchor=south] at (0.8,0.15) {Cl};
  \draw[->, thick, omDef] (1.6,1.75) to[bend right=40] node[midway, above] {e$^-$} (0,1.75);
\end{tikzpicture}

{\small 陶布的桥联中间体（示意）：仍与钴(III) 结合的氯离子，
通过取代一个水分子，进入铬(II) 的配位层，
电子经桥联传递，氯离子则留在
已成惰性的铬(III) 上。}
\end{omfigure}

电子跃迁约需 $10^{-16}$ 秒，远快于原子核的运动，
由弗兰克–康登原理（\cref{ch:b3:electronic-spectroscopy}）可知，
电子转移只能发生在核构型使反应物与生成物
具有相同能量的情形下。对于铁的自交换，\ce{Fe^{III}}--O 键短于
\ce{Fe^{II}}--O 键；在电子转移之前，两个配合物必须
畸变至共同的中间几何构型，并为此付出能量代价。

\section{马库斯理论}

\begin{definition}[重组能]\label{def:b3:complex-mechanisms:reorganisation}
电子转移的\emph{重组能}\index{重组能} $\lambda$ 是指：在不转移电子的情况下，
使反应物达到生成物的核构型（配合物的键长
与周围溶剂的取向）所需的能量。
它等于来自键的内层部分
与来自溶剂的外层部分之和。
\end{definition}

\begin{theorem}[马库斯方程]\label{thm:b3:complex-mechanisms:marcus}
若反应物与生成物的吉布斯自由能是曲率相同的抛物线，
沿反应坐标 $x$（反应物平衡处为 0，生成物平衡处为 1），
则标准反应吉布斯自由能为 $\Delta G^\circ$、
重组能为 $\lambda$ 的电子转移，其活化吉布斯自由能为
\[
  \Delta G^{\ddagger} = \frac{(\lambda + \Delta G^\circ)^2}{4\lambda} .
\]
此即\emph{马库斯方程}\index{马库斯方程}。
\end{theorem}

\begin{proof}
写 $G_R = \lambda x^2$，$G_P = \lambda(x - 1)^2 + \Delta G^\circ$：曲率由 $G_R(1) = \lambda$ 确定，即 $\lambda$ 的定义。
电子转移发生在两条曲线相交处：$\lambda x^2 = \lambda x^2 - 2\lambda x + \lambda +
\Delta G^\circ$，故 $x^\ast = (\lambda + \Delta G^\circ)/2\lambda$，且 $\Delta G^{\ddagger} = G_R(x^\ast) = (\lambda + \Delta G^\circ)^2/4\lambda$。
\end{proof}

\begin{corollary}[反转区]\label{cor:b3:complex-mechanisms:inverted}
在 $\lambda$ 固定时，反应越放能，速率增大，直至
$-\Delta G^\circ = \lambda$，此时 $\Delta G^{\ddagger} = 0$；超过此点后，速率又减小。
\end{corollary}

\begin{proof}
$\Delta G^{\ddagger}$ 是关于 $\Delta G^\circ$ 的抛物线，其最小值为零，在 $\Delta G^\circ = -\lambda$ 处取得；当 $\Delta G^\circ < -\lambda$ 时
它又增大。
\end{proof}

\begin{definition}[反转区]\label{def:b3:complex-mechanisms:inverted}
电子转移的\emph{反转区}\index{反转区}是指
驱动力 $-\Delta G^\circ > \lambda$ 的范围，在该范围内
反应越放能，速率反而减小。
\end{definition}

\begin{omfigure}
\begin{tikzpicture}
\begin{axis}[name=mp, width=0.47\linewidth, height=5.6cm, xmin=-0.6, xmax=2.2, ymin=-160, ymax=120,
  axis lines=left, xlabel={反应坐标 $x$}, ylabel={$G$ / \unit{kJ.mol^{-1}}},
  tick label style={font=\small}, label style={font=\small}, xtick={0,1,2}]
\addplot[black, thick] table[x=x, y=GR] {figdata/out/bachelor-3-marcus-parabolas.dat};
\addplot[omDef!50, thick] table[x=x, y=P0] {figdata/out/bachelor-3-marcus-parabolas.dat};
\addplot[omDef, thick] table[x=x, y=P50] {figdata/out/bachelor-3-marcus-parabolas.dat};
\addplot[omThm, thick] table[x=x, y=P100] {figdata/out/bachelor-3-marcus-parabolas.dat};
\addplot[omThm!60, thick, dashed] table[x=x, y=P150] {figdata/out/bachelor-3-marcus-parabolas.dat};
\node[font=\scriptsize, anchor=north] at (axis cs:0,-2) {反应物};
\end{axis}
\begin{axis}[at={(mp.east)}, anchor=west, xshift=1.4cm, width=0.47\linewidth, height=5.6cm,
  xmin=0, xmax=250, ymin=0, ymax=11.5, axis lines=left, xlabel={$-\Delta G^\circ$ / \unit{kJ.mol^{-1}}},
  ylabel={$\log(k/\unit{s^{-1}})$}, tick label style={font=\small}, label style={font=\small}]
\addplot[omDef, very thick] table[x=mdG, y=logk] {figdata/out/bachelor-3-marcus-rate.dat};
\draw[black!40, dashed] (axis cs:100,0) -- (axis cs:100,11);
\node[font=\scriptsize, anchor=south east] at (axis cs:95,1) {正常区};
\node[font=\scriptsize, anchor=south west] at (axis cs:105,1) {反转区};
\end{axis}
\end{tikzpicture}

{\small 马库斯理论，取 $\lambda = \qty{100}{kJ/mol}$（模型）。左：反应物抛物线
（黑色）以及 $\Delta G^\circ = 0$ 和 $-50$ 对应的生成物抛物线（蓝色）、$-100 = -\lambda$（红色）和
\qty{-150}{kJ/mol}（虚线，反转区）；它们的交点
即过渡态，其位置降至反应物曲线的最低点，
在 $-\Delta G^\circ = \lambda$ 处，之后又升高。右：速率常数的对数
（模型指前因子 \qty{e11}{s^{-1}}）在 $-\Delta G^\circ = \lambda$ 处达到极大值。}
\end{omfigure}

\begin{proposition}[外层重组能]\label{prop:b3:complex-mechanisms:outer-reorganisation}
对于两个球形反应物，半径分别为 $a_1$、$a_2$，二者相距 $d$，
置于折射率为 $n$、相对介电常数为 $\varepsilon_s$ 的溶剂中，$\lambda$ 的溶剂部分
与 $\bigl(\frac{1}{2a_1} + \frac{1}{2a_2} - \frac1d\bigr)\bigl(\frac{1}{n^2} - \frac{1}{\varepsilon_s}\bigr)$ 成正比。
\end{proposition}

\admitted

水的极性很强（$\varepsilon_s$ 很大），但折射率普通，故溶剂重组能较大；
非极性溶剂则较小。大的反应物重组得较少：
这正是电子转移蛋白把金属位点埋入
蛋白质内部的原因。

\begin{theorem}[马库斯交叉关系]\label{thm:b3:complex-mechanisms:cross-relation}
对于交叉反应 $\mathrm{A_{ox} + B_{red} \to A_{red} + B_{ox}}$，平衡常数为 $K_{12}$，两电对
的自交换速率常数分别为 $k_{11}$ 和 $k_{22}$，
\[
  k_{12} = \sqrt{k_{11}k_{22}K_{12}f}, \qquad \ln f = \frac{(\ln K_{12})^2}{4\ln(k_{11}k_{22}/Z^2)},
\]
其中 $Z$ 为碰撞频率因子；$f \approx 1$，前提是 $K_{12}$ 不太大。此即
\emph{马库斯交叉关系}\index{马库斯交叉关系}。
\end{theorem}

\begin{proof}[部分证明]
假设交叉反应的重组能为两个
自交换反应重组能的平均值，即 $\lambda_{12} = (\lambda_{11} + \lambda_{22})/2$（每个反应物贡献其自身的一半），
且所有速率常数均取 $k = Z\eu^{-\Delta G^{\ddagger}/RT}$。于是
$\Delta G_{12}^{\ddagger} = \frac{\lambda_{12}}{4}\bigl(1 + \frac{\Delta G^\circ_{12}}{\lambda_{12}}\bigr)^2 = \frac{\lambda_{12}}{4} + \frac{\Delta G^\circ_{12}}{2} + \frac{(\Delta G^\circ_{12})^2}{4\lambda_{12}}$。第一项给出
$\sqrt{k_{11}k_{22}}$（因 $\Delta G^{\ddagger}_{ii} = \lambda_{ii}/4$），第二项为 $\sqrt{K_{12}}$（取 $\Delta G^\circ_{12} = -RT\ln K_{12}$），第三项为
因子 $f$。
\end{proof}

\begin{method}[由自交换数据求交叉反应速率]\label{met:b3:complex-mechanisms:marcus-estimate}
\begin{enumerate}
  \item 求两个电对的自交换速率常数 $k_{11}$、$k_{22}$。
  \item 由标准电极电势计算 $K_{12}$：$\ln K_{12} = nF(E^\circ_{\text{氧化剂}} - E^\circ_{\text{还原剂}})/RT$。
  \item $k_{12} = \sqrt{k_{11}k_{22}K_{12}}$（乘以 $f$ 仅在 $K_{12}$ 很大时才需要）。与实测值相差
        在几倍以内，即表明为外层机理。
\end{enumerate}
\end{method}

马库斯于 1956 年发表了该理论；反转区起初受到质疑，
1984 年，格哈德·克洛斯与约翰·米勒在给体与受体
由刚性间隔基固定于一定距离的分子中观察到了它，
扩散便无法掩盖它。

\begin{inthelab}[氧-17 核磁共振测定水交换]
与顺磁性离子配位的水的氧-17 共振信号，以及本体水的共振信号，
因两者之间的交换而展宽。记录本体水信号的线宽
随温度的变化，所用溶液为高氯酸金属盐，
其中的水以 \ce{^{17}O} 富集，由此可得交换速率常数；若用高压探头，
还可测得其活化体积。
\end{inthelab}

\begin{safety}{\ghs[0.8cm]{GHS05}\ghs[0.8cm]{GHS06}\ghs[0.8cm]{GHS08}}
% ledger: ghs4:K2PtCl4
四氯合铂（II）酸钾：吞食有毒，可引起严重的
眼损伤，并为呼吸道和皮肤致敏物。铂配合物须在通风橱中
戴手套操作；顺铂及其类似物为细胞毒性药物，须按此类药物
处理。
\end{safety}

\begin{history}[陶布与马库斯]
亨利·陶布于 1952 年将配合物分为活性与惰性两类，
并于 1953 年以铬–钴实验确立了内层机理；他于 1983 年
获诺贝尔化学奖。鲁道夫·马库斯于 1956 年推导出
以其名字命名的电子转移理论，并于 1992 年获得诺贝尔
化学奖。
\end{history}

\section{练习}

\begin{exercise}[$\star$]\label{exo:b3:complex-mechanisms:1}
判断下列配合物为活性还是惰性，并说明理由：$\ce{[Cr(H2O)6]^{3+}}$、$\ce{[Co(NH3)6]^{3+}}$、$\ce{[Cr(H2O)6]^{2+}}$，
$\ce{[Cu(H2O)6]^{2+}}$、$\ce{[Zn(H2O)6]^{2+}}$。
\end{exercise}

\begin{exercise}[$\star$]\label{exo:b3:complex-mechanisms:2}
证明：D 速率方程化为 $v = k_1[\mathrm{ML_5X}]$（$[\mathrm Y]$ 很大时），并化为 $v = (k_1k_2/k_{-1})
[\mathrm{ML_5X}][\mathrm Y]/[\mathrm X]$（当 $k_{-1}[\mathrm X] \gg k_2[\mathrm Y]$ 时）。
\end{exercise}

\begin{exercise}[$\star$]\label{exo:b3:complex-mechanisms:3}
预测 D 型与 A 型取代的 $\Delta V^{\ddagger}$ 与 $\Delta^{\ddagger}S^\circ$ 的符号。
\end{exercise}

\begin{exercise}[$\star$]\label{exo:b3:complex-mechanisms:4}
写出 $\ce{[PtCl4]^{2-}}$ 与两个 \ce{NH3} 反应的产物，以及 $\ce{[Pt(NH3)4]^{2+}}$ 与两个 \ce{Cl-} 反应的产物。
\end{exercise}

\begin{exercise}[$\star\star$]\label{exo:b3:complex-mechanisms:5}
对水合离子上的取代，$K_{\mathrm{os}} = \qty{2.0}{L.mol^{-1}}$，$k_i = \qty{3.0e4}{s^{-1}}$（
本题数据）。求 $k_{\mathrm{obs}}$，取 $[\mathrm Y] = 0.10$ 与 \qty{1.0}{mol/L}，并求其极限。
\end{exercise}

\begin{exercise}[$\star\star$]\label{exo:b3:complex-mechanisms:6}
在 \qty{298}{K} 下，压强由 0.1 升至 \qty{100}{MPa} 时，某速率常数增大一倍。
求 $\Delta V^{\ddagger}$，并作出结论。
\end{exercise}

\begin{exercise}[$\star\star$]\label{exo:b3:complex-mechanisms:7}
在 $\ce{[PtCl3(PEt3)]-}$ 中，与 \ce{PEt3} 呈反位的 Pt--Cl 键为 \qty{2.42}{\angstrom}，其余两个为 \qty{2.32}{\angstrom}（
本题数据）。试解释此现象，并预测哪一个氯配体先被取代。
\end{exercise}

\begin{exercise}[$\star\star$]\label{exo:b3:complex-mechanisms:8}
列出能表明某一电子转移属于内层
机理的实验证据。
\end{exercise}

\begin{exercise}[$\star\star$]\label{exo:b3:complex-mechanisms:9}
取 $\lambda = \qty{80}{kJ/mol}$，求 $\Delta G^{\ddagger}$，分别对应 $\Delta G^\circ = 0$、$-20$、$-80$ 及 \qty{-120}{kJ/mol}。
\end{exercise}

\begin{exercise}[$\star\star\star$]\label{exo:b3:complex-mechanisms:10}
两个电对的自交换速率常数分别为 $k_{11} = 4.0$ 与 $k_{22} = \qty{2.0e5}{L.mol^{-1}.s^{-1}}$，且
其交叉反应的 $K_{12} = \num{1.0e4}$（本题数据）。估算 $k_{12}$（取 $f = 1$）。
\end{exercise}

\begin{exercise}[$\star\star\star$]\label{exo:b3:complex-mechanisms:11}
采用仅考虑 $\sigma$ 的模型，计算
高自旋 $d^5$、$d^7$ 与 $d^8$ 离子的配体场活化能，并按预期水交换速率对
\ce{Mn^{2+}}、\ce{Co^{2+}} 与 \ce{Ni^{2+}} 排序。
\end{exercise}

\begin{exercise}[$\star\star\star$]\label{exo:b3:complex-mechanisms:12}
为何溶液中相互扩散的离子间反应难以观察到反转区，
而刚性给体–受体分子
又是如何将其揭示出来的？
\end{exercise}

\section{周末问题：制备顺铂而非反铂}

\begin{problem}[{周末问题——制备顺铂而非反铂：平面正方形的
铂（II）离子、借助反位效应设计的合成、药物的水合作用，
以及它在细胞内被活化的原因}]
\label{pb:b3:complex-mechanisms:1}
% ledger: ghs4:K2PtCl4
顺铂 $cis$-$\ce{[PtCl2(NH3)2]}$ 由 $\ce{K2[PtCl4]}$ 制得。本题数据：
在 \qty{37}{\celsius} 时，第一步水合反应 $\ce{[PtCl2(NH3)2] + H2O -> [PtCl(H2O)(NH3)2]+ + Cl-}$ 的
速率常数为 $k = \qty{9.0e-5}{s^{-1}}$，平衡常数为 $K = \qty{3.0e-3}{mol/L}$；
血浆中氯离子浓度约为 \qty{100}{mM}，细胞内约为 \qty{4}{mM}。

\medskip\noindent\textbf{第一部分——铂（II）。}
\begin{enumerate}
  \item 写出 \ce{Pt^{2+}} 的 $d$ 电子数。
  \item 为何其配合物呈平面正方形且为抗磁性？
  \item 为何其取代反应为缔合型？
  \item 写出双项速率方程，并解释其中两项的含义。
  \item 以一天为时间尺度，该药物是活性的还是惰性的？
\end{enumerate}

\medskip\noindent\textbf{第二部分——合成。}
\begin{enumerate}[resume]
  \item $\ce{K2[PtCl4]}$ 与两个 \ce{NH3} 反应会得到什么？为什么？
  \item 这种直接路线得到的产物不纯。达拉的路线先用过量 KI，
        将 $\ce{[PtCl4]^{2-}}$ 转化为 $\ce{[PtI4]^{2-}}$。为何碘离子有用？
  \item 加入 \ce{NH3} 得到 $cis$-$\ce{[PtI2(NH3)2]}$。请解释其立体化学。
  \item 随后用硝酸银水溶液移除碘化物。生成了什么？
  \item 加入 KCl 得到顺铂。为什么构型不变？
  \item 那么，应如何制得反铂？
  \item 反铂也能发生反应，为何不能用作药物？
\end{enumerate}

\medskip\noindent\textbf{第三部分——水合作用。}
\begin{enumerate}[resume]
  \item 计算第一步水合反应的半衰期。
  \item 计算在无氯离子溶液中经过 \qty{2.0}{h} 后被水合的分数。
  \item 为何水合配合物是与 DNA 反应的活性形式？
  \item 该水合反应是解离型还是缔合型？其
        活化体积应为多少？
  \item 注射用药物溶液为何用生理盐水配制？
\end{enumerate}

\medskip\noindent\textbf{第四部分——血浆与细胞。}
\begin{enumerate}[resume]
  \item 将水合配合物的平衡分数写成
        $[\ce{Cl-}]$ 的函数。
  \item 计算其在血浆中的值。
  \item 计算其在细胞内的值。
  \item 计算二者之比。
  \item 因此，该药物为何主要在细胞内发挥作用？
  \item 细胞内还有哪些配体能与铂结合，
        并使药物失活？
  \item 为何两个氨配体始终保持结合？
  \item 陈述结论：细胞内水合配合物的分数与
        血浆中的分数之比。
\end{enumerate}
\end{problem}
